% Student-only reader-number cross-references for public version 1.0.
\DeclarePublicXref{CM-C01-U001}{%
%
\begin{PublicXrefBox}{Para practicar}
\textbf{Ejercicios recomendados:} 1.1, 1.3 y 1.2.\par
\textbf{Colección práctica:} \texttt{P01}.\par
Clasifica, calcula y justifica antes de avanzar a la recta real.
\end{PublicXrefBox}%

}
\DeclarePublicXref{CM-C01-U007}{%
%
\begin{PublicXrefBox}{Para practicar}
\textbf{Ejercicios recomendados:} 1.14, 1.17 y 1.27.\par
\textbf{Colección práctica:} \texttt{P02}.\par
Trabaja distancia, cotas e imagen/preimagen con una justificación visible.
\end{PublicXrefBox}%

}
\DeclarePublicXref{CM-C01-U012}{%
%
\begin{PublicXrefBox}{Para practicar}
\textbf{Ejercicios recomendados:} 1.31, 1.34 y 1.41.\par
\textbf{Colección práctica:} \texttt{P03}.\par
Declara dominios y codominios antes de invertir, componer o resolver.
\end{PublicXrefBox}%

}
\DeclarePublicXref{CM-C02-U003}{%
%
\begin{PublicXrefBox}{Para practicar}
\textbf{Ejercicios recomendados:} 2.3, 2.9 y 2.8.\par
\textbf{Colección práctica:} \texttt{P04}.\par
Pasa de una descripción de sucesión a un cálculo y después a un argumento epsilon--N.
\end{PublicXrefBox}%

}
\DeclarePublicXref{CM-C02-U011}{%
%
\begin{PublicXrefBox}{Para practicar}
\textbf{Ejercicios recomendados:} 2.19, 2.30 y 2.37.\par
\textbf{Colección práctica:} \texttt{P05}.\par
Alterna una propiedad estructural, una prueba con delta y una lectura de límites laterales.
\end{PublicXrefBox}%

}
\DeclarePublicXref{CM-C02-U016}{%
%
\begin{PublicXrefBox}{Para practicar}
\textbf{Ejercicios recomendados:} 2.45, 2.53 y 2.58.\par
\textbf{Colección práctica:} \texttt{P06}.\par
Transforma el límite, verifica la composición y clasifica la discontinuidad.
\end{PublicXrefBox}%

}
\DeclarePublicXref{CM-C02-U019}{%
%
\begin{PublicXrefBox}{Para practicar}
\textbf{Ejercicios recomendados:} 2.61, 2.64 y 2.68.\par
\textbf{Material de consolidación:} \texttt{P07} (opcional; no programado como sesión docente en 2026/27).\par
Comprueba hipótesis y distingue existencia, localización y unicidad.
\end{PublicXrefBox}%

}
\DeclarePublicXref{CM-C03-U001}{%
%
\begin{PublicXrefBox}{Para practicar}
\textbf{Ejercicios recomendados:} 3.4, 3.2 y 3.1.\par
\textbf{Colección práctica:} \texttt{P08}.\par
Reconstruye una derivada y conecta pendiente, tasa y significado local.
\end{PublicXrefBox}%

}
\DeclarePublicXref{CM-C03-U006}{%
%
\begin{PublicXrefBox}{Para practicar}
\textbf{Ejercicios recomendados:} 3.7, 3.11 y 3.19.\par
\textbf{Colección práctica:} \texttt{P09}.\par
Combina linealización, relación entre derivabilidad y continuidad, y reglas algebraicas.
\end{PublicXrefBox}%

}
\DeclarePublicXref{CM-C03-U012}{%
%
\begin{PublicXrefBox}{Para practicar}
\textbf{Ejercicios recomendados:} 3.25, 3.31 y 3.32.\par
\textbf{Colección práctica:} \texttt{P10}.\par
Identifica la estructura antes de elegir cadena, inversa o derivación implícita.
\end{PublicXrefBox}%

}
\DeclarePublicXref{CM-C04-U005}{%
%
\begin{PublicXrefBox}{Para practicar}
\textbf{Ejercicios recomendados:} 3.49, 3.52 y 4.10.\par
\textbf{Colección práctica:} \texttt{P11}.\par
Verifica las hipótesis del valor medio y de l'Hôpital antes de estudiar variación.
\end{PublicXrefBox}%

}
\DeclarePublicXref{CM-C04-U014}{%
%
\begin{PublicXrefBox}{Para practicar}
\textbf{Ejercicios recomendados:} 4.26, 4.32 y 4.42.\par
Clasifica extremos, compara candidatos y conecta la segunda derivada con convexidad.
\end{PublicXrefBox}%

}
\DeclarePublicXref{CM-C05-U007}{%
%
\begin{PublicXrefBox}{Para practicar}
\textbf{Ejercicios recomendados:} 5.3, 5.19 y 5.4.\par
\textbf{Colección práctica:} \texttt{P12}.\par
Predice el estudio cualitativo, justifica las asíntotas y contrasta al final con una herramienta.
\end{PublicXrefBox}%

}
\DeclarePublicXref{CM-C06-U003}{%
%
\begin{PublicXrefBox}{Para practicar}
\textbf{Ejercicios recomendados:} 6.3, 6.5 y 6.10.\par
Relaciona orden de contacto, construcción por derivadas y polinomios básicos.
\end{PublicXrefBox}%

}
\DeclarePublicXref{CM-C06-U012}{%
%
\begin{PublicXrefBox}{Para practicar}
\textbf{Ejercicios recomendados:} 6.13, 6.19, 6.27 y 6.32.\par
\textbf{Colección práctica:} \texttt{P13}.\par
Comprueba las hipótesis, elige el grado desde una cota y conserva el primer término no nulo.
\end{PublicXrefBox}%

}
